\section{Prepare a cumulative observation}
Start from \texttt{packages/023-delta-varint/starter/main.asm}. Replace the named region's
temporary two lines with the corresponding increment, keeping labels and markers.
The frames and status exits are supplied. One special initial width gate branches
forward to \texttt{dec\_accumulate}; it deliberately bypasses fifth-byte validation
until the payload loop is installed. Use only the specified known-valid recipes at
that stage. Replace that gate when its instruction step arrives.

Keep the cumulative source's count, capacity and method declarations at their defaults.
Use \texttt{prepare.py} with a recipe ID, your cumulative source path and a separate
output path. It writes the exact logical input and selects the proper function for
that observation without overwriting your source or an existing file. Input octets
are listed unsigned in the recipes; the prepared assembly spells high byte values
as negative byte initializers with the same low eight bits. The decoder recipes fill
the word destination with a sentinel and the encoder recipes fill the byte destination
with \texttt{a5}.

Most observations stop at \texttt{codec\_done}, after the supplied restoring exit
and before caller result stores. The recipe gives returned counts as input units,
output units and status. Temporary registers retain the illustrated partial state;
saved registers have already been restored. Recipe \texttt{s57-1} instead stops at
\texttt{dec\_terminal} inside the completed decoder and inspects current working
\texttt{s5}, \texttt{s6} before restoration. Continue to \texttt{results\_ready} to
see the caller's three result words populated.

A temporary exit is a construction aid. Some emission/store steps precede their
counter updates, so a partial program can contain pending output bytes while returning
zero or mismatched counters. Those observations describe only that named stage.
They do not authorize partial publication by the completed codec. Final behavioral
checks independently verify item-atomic stores, counts, status priorities and all
preservation requirements.

The illustrated allocation uses \texttt{s0} for the full source base, \texttt{s1} for
input units, \texttt{s2} for the full destination base, \texttt{s3} for output capacity,
\texttt{s4} for the committed predecessor, \texttt{s5} for committed input and
\texttt{s6} for committed output. Encoder \texttt{t0} retains the candidate reading;
decoder \texttt{t0} is its private byte scan, \texttt{t1} the unsigned accumulator,
\texttt{t2} the payload shift and \texttt{t3} its candidate length. Equivalent completed
solutions can use other internal labels and register allocations.
